\section{Contoh Trace Lengkap AES-128 (Satu Round)}

\subsection{Setup Contoh}

Untuk memahami AES secara mendalam, kita akan melakukan trace lengkap satu round AES-128 dengan nilai konkret. Ini akan menunjukkan bagaimana setiap operasi (SubBytes, ShiftRows, MixColumns, AddRoundKey) mengubah state.

\textbf{Input}:
\begin{itemize}
  \item Plaintext: \texttt{00 11 22 33 44 55 66 77 88 99 AA BB CC DD EE FF}
  \item Key: \texttt{00 01 02 03 04 05 06 07 08 09 0A 0B 0C 0D 0E 0F}
\end{itemize}

\subsection{State Awal (Setelah AddRoundKey Round 0)}

State direpresentasikan sebagai matriks 4x4 (kolom-major order):

\begin{table}[htbp]
  \centering
  \small
  \begin{tabular}{|c|c|c|c|}
    \hline
    \textbf{00} & \textbf{44} & \textbf{88} & \textbf{CC} \\
    \hline
    \textbf{11} & \textbf{55} & \textbf{99} & \textbf{DD} \\
    \hline
    \textbf{22} & \textbf{66} & \textbf{AA} & \textbf{EE} \\
    \hline
    \textbf{33} & \textbf{77} & \textbf{BB} & \textbf{FF} \\
    \hline
  \end{tabular}
  \caption{State awal setelah AddRoundKey dengan round key 0.}
\end{table}

\subsection{Step 1: SubBytes}

Setiap byte disubstitusi menggunakan S-box AES. Beberapa contoh substitusi:
\begin{itemize}
  \item \texttt{00} $\rightarrow$ \texttt{63} (S-box[0x00] = 0x63)
  \item \texttt{11} $\rightarrow$ \texttt{82} (S-box[0x11] = 0x82)
  \item \texttt{22} $\rightarrow$ \texttt{C9} (S-box[0x22] = 0xC9)
  \item \texttt{33} $\rightarrow$ \texttt{7B} (S-box[0x33] = 0x7B)
\end{itemize}

\begin{table}[htbp]
  \centering
  \small
  \begin{tabular}{|c|c|c|c|}
    \hline
    \textbf{63} & \textbf{09} & \textbf{CD} & \textbf{BA} \\
    \hline
    \textbf{82} & \textbf{74} & \textbf{60} & \textbf{3A} \\
    \hline
    \textbf{C9} & \textbf{E9} & \textbf{D0} & \textbf{E0} \\
    \hline
    \textbf{7B} & \textbf{C3} & \textbf{66} & \textbf{16} \\
    \hline
  \end{tabular}
  \caption{State setelah SubBytes.}
\end{table}

\subsection{Step 2: ShiftRows}

Setiap baris digeser ke kiri dengan jumlah berbeda:
\begin{itemize}
  \item Baris 0: tidak digeser
  \item Baris 1: geser 1 posisi $\rightarrow$ [82, 74, 60, 3A] menjadi [74, 60, 3A, 82]
  \item Baris 2: geser 2 posisi $\rightarrow$ [C9, E9, D0, E0] menjadi [D0, E0, C9, E9]
  \item Baris 3: geser 3 posisi $\rightarrow$ [7B, C3, 66, 16] menjadi [16, 7B, C3, 66]
\end{itemize}

\begin{table}[htbp]
  \centering
  \small
  \begin{tabular}{|c|c|c|c|}
    \hline
    \textbf{63} & \textbf{09} & \textbf{CD} & \textbf{BA} \\
    \hline
    \textbf{74} & \textbf{60} & \textbf{3A} & \textbf{82} \\
    \hline
    \textbf{D0} & \textbf{E0} & \textbf{C9} & \textbf{E9} \\
    \hline
    \textbf{16} & \textbf{7B} & \textbf{C3} & \textbf{66} \\
    \hline
  \end{tabular}
  \caption{State setelah ShiftRows.}
\end{table}

\subsection{Step 3: MixColumns}

Setiap kolom dikalikan dengan matriks MixColumns dalam GF($2^8$). Untuk kolom pertama [63, 74, D0, 16]:

\[
\begin{bmatrix}
02 & 03 & 01 & 01 \\
01 & 02 & 03 & 01 \\
01 & 01 & 02 & 03 \\
03 & 01 & 01 & 02
\end{bmatrix}
\cdot
\begin{bmatrix}
63 \\ 74 \\ D0 \\ 16
\end{bmatrix}
=
\begin{bmatrix}
5F \\ 72 \\ 64 \\ 15
\end{bmatrix}
\]

\textbf{Perhitungan detail untuk elemen pertama}:
\begin{align*}
s'_{0,0} &= (02 \cdot 63) \oplus (03 \cdot 74) \oplus (01 \cdot D0) \oplus (01 \cdot 16) \\
&= C6 \oplus 9C \oplus D0 \oplus 16 \\
&= 5F
\end{align*}

dimana perkalian dalam GF($2^8$) menggunakan polynomial irreducible $x^8 + x^4 + x^3 + x + 1$ (0x11B).

\begin{table}[htbp]
  \centering
  \small
  \begin{tabular}{|c|c|c|c|}
    \hline
    \textbf{5F} & \textbf{9F} & \textbf{16} & \textbf{38} \\
    \hline
    \textbf{72} & \textbf{E6} & \textbf{5D} & \textbf{D4} \\
    \hline
    \textbf{64} & \textbf{AF} & \textbf{4E} & \textbf{A4} \\
    \hline
    \textbf{15} & \textbf{22} & \textbf{8B} & \textbf{2C} \\
    \hline
  \end{tabular}
  \caption{State setelah MixColumns.}
\end{table}

\subsection{Step 4: AddRoundKey}

State di-XOR dengan round key 1. Misalkan round key 1 adalah:

\begin{table}[htbp]
  \centering
  \small
  \begin{tabular}{|c|c|c|c|}
    \hline
    \textbf{D6} & \textbf{AA} & \textbf{74} & \textbf{FD} \\
    \hline
    \textbf{D2} & \textbf{63} & \textbf{9C} & \textbf{CA} \\
    \hline
    \textbf{E0} & \textbf{97} & \textbf{43} & \textbf{B4} \\
    \hline
    \textbf{04} & \textbf{02} & \textbf{38} & \textbf{51} \\
    \hline
  \end{tabular}
  \caption{Round key 1 (dari key expansion).}
\end{table}

XOR state dengan round key:

\begin{align*}
\text{State}[0][0] &= 5F \oplus D6 = 89 \\
\text{State}[1][0] &= 72 \oplus D2 = A0 \\
\text{State}[2][0] &= 64 \oplus E0 = 84 \\
\text{State}[3][0] &= 15 \oplus 04 = 11
\end{align*}

\begin{table}[htbp]
  \centering
  \small
  \begin{tabular}{|c|c|c|c|}
    \hline
    \textbf{89} & \textbf{35} & \textbf{62} & \textbf{C5} \\
    \hline
    \textbf{A0} & \textbf{85} & \textbf{C1} & \textbf{1E} \\
    \hline
    \textbf{84} & \textbf{38} & \textbf{0D} & \textbf{10} \\
    \hline
    \textbf{11} & \textbf{20} & \textbf{B3} & \textbf{7D} \\
    \hline
  \end{tabular}
  \caption{State setelah AddRoundKey (akhir round 1).}
\end{table}

\subsection{Ringkasan Transformasi}

\begin{table}[htbp]
  \centering
  \small
  \begin{tabular}{lp{8cm}}
    \toprule
    \textbf{Operasi} & \textbf{Tujuan} \\
    \midrule
    SubBytes & Substitusi non-linear untuk confusion \\
    ShiftRows & Permutasi untuk diffusion antar kolom \\
    MixColumns & Mixing linear dalam GF($2^8$) untuk diffusion dalam kolom \\
    AddRoundKey & Kombinasi dengan kunci untuk key-dependent transformation \\
    \bottomrule
  \end{tabular}
  \caption{Ringkasan tujuan setiap operasi AES.}
\end{table}

\begin{catatan}
Trace lengkap ini menunjukkan bagaimana AES mengubah plaintext melalui kombinasi operasi non-linear (SubBytes), permutasi (ShiftRows), mixing (MixColumns), dan key addition (AddRoundKey). Setiap operasi berkontribusi pada confusion dan diffusion yang membuat AES sangat aman.
\end{catatan}
